\section{Introduction}\label{sec:intro}

The exponential growth of digital video archives necessitates fine-grained, semantically robust multimodal retrieval mechanisms. Traditional metadata-based search paradigms are fundamentally inadequate for capturing the rich spatial and temporal dynamics inherent in modern video collections. Consequently, the field has increasingly pivoted toward the Known-Item Search (KIS) paradigm and Video Question Answering, demanding systems capable of pinpointing highly specific video segments from vast, unstructured repositories based on complex natural language queries. While vision-language models (VLMs) \cite{radford2021learning, zhai2023sigmoid} offer unprecedented zero-shot retrieval capabilities, deploying these architectures in real-world, linguistically diverse, and structurally rigorous environments introduces critical bottlenecks.

This paper addresses four fundamental research challenges in contemporary multimodal retrieval systems. The first challenge, \textit{Research Challenge 1 (RC1): The Cross-Lingual Semantic Gap}, arises when users submit queries in low-resource languages or dialects, such as Vietnamese, frequently characterized by missing diacritics and typographical errors. Conversely, state-of-the-art VLMs are predominantly pretrained on massive English corpora, including LAION-5B \cite{schuhmann2022laion} and WebLI. Direct lexical translation often corrupts domain-specific nuances, resulting in severe semantic degradation during the query embedding phase.

\textit{Research Challenge 2 (RC2): Generative Hallucination in Query Expansion}. While techniques like Hypothetical Document Embeddings (HyDE) \cite{gao2023precise} and Large Language Model (LLM) based query expansion have proven effective in dense retrieval \cite{chen2024bge}, they are notoriously prone to fabricating concrete details such as named entities, license plates, and specific geographical locations. Such unconstrained, speculative expansion directly contaminates the embedding space, misleading the retrieval engine away from factual ground truth.

\textit{Research Challenge 3 (RC3): Temporal Autocorrelation and Visual Redundancy}. Continuous video streams exhibit profound temporal autocorrelation. Consequently, unfiltered Approximate Nearest Neighbor (ANN) search frequently floods the top-$K$ retrieval results with near-identical frames extracted from the same static scene, drastically reducing recall for diverse candidate events. Furthermore, complex, multi-step temporal queries (e.g., event A followed by event B, preceding event C) lack native operational semantics and execution structures in standard vector database architectures.

\textit{Research Challenge 4 (RC4): Auditability and the Provenance Void}. Modern neural vector search pipelines operate predominantly as opaque transformations, systematically discarding critical intermediate metadata during the embedding process. Extracted intelligence such as Optical Character Recognition (OCR) bounding boxes, Automatic Speech Recognition (ASR) timestamps, and source file checksums are decoupled from the final vector representations. This architectural flaw renders robust error diagnosis and deterministic provenance tracking virtually impossible.

To address these systemic deficiencies, we introduce Vecna, a comprehensive multimodal video retrieval framework whose end-to-end architecture and operational pipeline are outlined in Figure~\ref{fig:architecture}. The core contributions of this work are summarized as follows:
\begin{itemize}
    \item We propose a Multi-Aspect Factual HyDE methodology, strictly governed by a Zero-Speculation Principle, which deterministically generates orthogonal 5-tuple query representations to mitigate generative hallucination.
    \item We formulate a two-pointer sliding-window Dynamic Programming (DP) algorithm for precise temporal sequence alignment, achieving an optimal $\mathcal{O}(|C|+|B|)$ amortized computational complexity.
    \item We develop a quantile cosine-smoothing scene segmentation algorithm paired with segment-level diversification strategies to counteract temporal autocorrelation and maximize top-$K$ variance.
    \item We implement an end-to-end cryptographic provenance sidecar architecture, ensuring SHA-256 fixity and immutable OCR/ASR evidence binding throughout the retrieval pipeline.
\end{itemize}

The remainder of this paper is organized as follows. Section \ref{sec:related} reviews related work in multimodal retrieval and vision-language alignment. Section \ref{sec:arch} describes the system architecture. Section \ref{sec:method} details the mathematical formulation of our proposed methodology. Section \ref{sec:exp} presents our rigorous experimental setup and empirical results. Finally, Section \ref{sec:conclusion} concludes the paper and discusses avenues for future research.
